\section{Background}
\label{sec:background}

We provide relevant background on computer security and AI agents.

\subsection{Computer Security}

In this work, we focus on the \emph{vulnerability exploitation} of computer systems.
A \emph{vulnerability} in a computer system is flaw in that system that allows
behaviors unintended by the creator of the system, typically for 
malicious use. \emph{Exploiting} the vulnerability consists of \emph{detecting} the
vulnerability and performing the necessary actions to take advantage of the
vulnerability. 

We focus on vulnerabilities in a computer system that are unknown to the
deployer of the system. Unfortunately, the term of these
vulnerabilities vary from source to source, but we refer to these
vulnerabilities as \emph{zero-day vulnerabilities} (0DV). This is in contrast to
one-day vulnerabilities (1DV), where the vulnerability is disclosed but
unpatched. Namely, a 1DV is \emph{known to the attacker}. 

Zero-day vulnerabilities are particularly harmful because the system deployer
cannot proactively put mitigations in place against these vulnerabilities
\cite{bilge2012before}. We focus specifically on web vulnerabilities in this
work, which are often the first attack surface into more in depth attacks
\cite{setiawan2018web}.

One important distinction within vulnerabilities is the \emph{class} of
vulnerability and the \emph{specific instance} of the vulnerability. For
example, server-side request forgery (SSRF) has been known as a class of
vulnerability since at least 2011 \cite{fung2011privacy}. However, one of the
biggest hacks of all time that occurred in 2021 (10 years after) hacked
Microsoft, now a multi-trillion dollar company that invests about a billion
dollars a year in computer security \cite{microsoft2024securing}, used an SSRF
\cite{kost2023critical}.

Thus, specific \emph{instances} of zero-day vulnerabilities are critical to
find.


\subsection{AI Agents and Cybersecurity}

AI agents have become increasingly powerful and can perform tasks as complex as
solving real-world GitHub issues \cite{yang2024sweagent}. In this work, we focus
on AI agents solving complex, real-world tasks. These agents are now almost
exclusively powered by tool-enabled LLMs \cite{parisi2022talm, weng2023agent}.
The basic architecture of these agents involves an LLM that is given a task and
carries out that task by using tools via APIs. We provide a more detailed
overview of AI agents in Section~\ref{sec:rel-work}.

Recent work has explored AI agents in the context of cybersecurity, showing that
they can exploit ``capture-the-flag'' style vulnerabilities \cite{fang2024llm,zhang2024cybench}
and one-day vulnerabilities when given a description of the vulnerability
\cite{fang2024llm2}. These agents work via the ReAct-style iteration,
where LLMs take an action, observe the response, and repeat \cite{yao2022react}.

However, these agents fare poorly in the zero-day setting. We now describe our
architecture for improving these agents.
